\chapter{El Motor ZINA 2084}

El sistema de juego de \textbf{NeoBilbao 2084} utiliza el motor narrativo emocional \textbf{ZINA}, una adaptación de las filosofías de diseño \emph{Powered by the Apocalypse} (PbtA) y \emph{The Veil}. En este sistema no existen atributos físicos convencionales como Fuerza, Agilidad o Inteligencia. Lo que determina el éxito o el fracaso de una acción es la \textbf{disposición emocional} del personaje ante el conflicto.

\section{Las seis emociones de 2084}

A diferencia de la ambientación de 2102 (donde la emoción «Alegre» está tachada para siempre), en 2084 la clase trabajadora aún conserva su capacidad para la esperanza, el compañerismo y el humor frente a la desgracia.

Los personajes reparten el conjunto de valores \textbf{(+2, +1, +1, 0, 0, -1)} entre las seis emociones:

\begin{itemize}
  \item \lead{Enfadado:} Violencia visceral, rebeldía, indignación y fuerza bruta. Se utiliza para embestir a guardias, romper cerraduras a patadas o iniciar peleas en la tasca.
  \item \lead{Asustado:} Instinto puro de supervivencia, reflejos de huida y prudencia. Se usa para escabullirse de patrullas de los Secretas, ocultarse tras contenedores o saltar entre tejados.
  \item \lead{Culpable:} El peso de las faltas pasadas, la mala conciencia y la necesidad de redención. Es el motor principal para vulnerar códigos corporativos, hackear terminales arriesgando la propia seguridad o confrontar la verdad.
  \item \lead{Afligido:} Empatía profunda, duelo compartido y dolor por los compañeros caídos. Se emplea para consolar a víctimas, tratar heridas traumáticas o convencer mediante la lágrima y la compasión.
  \item \lead{Pacífico:} Calma analítica, sangre fría, meticulosidad y desapego. Se utiliza para descifrar manuales técnicos, falsificar formularios burocráticos o pasar desapercibido en un control de seguridad.
  \item \lead{Alegre:} Esperanza activa, calor humano, camaradería de barrio y desafío moral. Se utiliza para infundir ánimos a un aliado al borde del colapso, mantener la cabeza alta frente a un interrogatorio o improvisar salidas ingeniosas.
\end{itemize}

\section{Resolución básica: 2d6 + Emoción}

Cuando un personaje realiza una acción arriesgada y la ficción exige una tirada, el Director determina qué emoción describe mejor el impulso del personaje. El jugador tira \textbf{2d6} y suma el modificador correspondiente:

\begin{CajaRegla}{La escala de resultados}
\begin{itemize}
  \item \textbf{10+ (Éxito total):} Consigues tu objetivo limpiamente. Tomas el control de la escena sin consecuencias negativas inesperadas.
  \item \textbf{7--9 (Éxito con coste):} Logras lo que pretendías, pero a un precio. Sufres daño físico, aumentas tu nivel de Zama, gastas una etiqueta de equipo valiosa o te colocas en una posición táctica desfavorable.
  \item \textbf{6 o menos (Fallo):} La situación empeora drásticamente. El Director realiza un movimiento duro: aparecen refuerzos corporativos, el sistema se bloquea o un compañero sufre las consecuencias.
\end{itemize}
\end{CajaRegla}

\section{Zama (Estrés) y Colapso}

La \textbf{Zama} representa la carga psicológica, neurotóxica y emocional que acumula el personaje. La escala va de 0 a 6:
\begin{itemize}
  \item Todo personaje comienza con \textbf{0 puntos de Zama}.
  \item El estrés aumenta por ataques neurotóxicos de los Auditores, visiones de la Resonancia Azul o fracasos traumáticos en la ficción.
  \item \lead{Al alcanzar 6 puntos de Zama:} El personaje sufre un \textbf{Colapso}. Queda fuera de combate: puede desmayarse en shock catatónico, sufrir un ataque de pánico incontrolable que delata la posición del grupo o acatar sumisamente las órdenes de un Auditor por miedo visceral.
  \item \lead{Recuperación de Zama:} La Zama se reduce descansando en lugares seguros (como la trastienda del bar «Zulo»), compartiendo kalimotxo y confidencias humanas, o mediante el uso de estimulantes médicos robados.
\end{itemize}

\section{La reserva de Auzolan (Solidaridad)}

El \textbf{Auzolan} es el antiguo principio vasco de ayuda vecinal comunitaria, adaptado a la resistencia urbana.
\begin{itemize}
  \item Al inicio de cada sesión de juego, el grupo recibe \textbf{1 punto de Auzolan por cada jugador presente}.
  \item Los puntos se colocan en el centro de la mesa en forma de fichas compartidas.
  \item Cualquier jugador puede gastar 1 punto de la reserva para:
    \begin{enumerate}
      \item \textbf{Prestar apoyo directo:} Sumar \textbf{+1} a la tirada de un compañero antes o después de lanzar los dados.
      \item \textbf{Declarar un hecho de barrio:} Declarar la existencia de un contacto barrial, un atajo por alcantarillas o una herramienta escondida en la escena.
      \item \textbf{Absorber el golpe:} Reducir en 1 punto el daño o la Zama que fuera a sufrir un aliado adyacente.
    \end{enumerate}
\end{itemize}

\section{Ejemplos de juego en La Margen}

\begin{CajaEjemplo}{Ejemplo 1: Violencia callejera frente a los Secretas}
\lead{Situación:} El Sabueso persigue a un Auditor por un callejón y este desenfunda un inyector neumático de amnésico.\\
\lead{Ficción:} El jugador describe cómo arremete contra el Auditor con una barra de hierro oxidada para desarmarlo por la fuerza.\\
\lead{Mecánica:} El Director pide tirar 2d6 + \textbf{Enfadado} (+2). El jugador saca un 6 en los dados, para un total de 8 (\emph{Éxito con coste}).\\
\lead{Resolución:} El Sabueso derriba al Auditor contra un contenedor de basura, pero en el forcejeo la aguja del inyector le roza el antebrazo. El Auditor queda inconsciente, pero el Sabueso sufre 1 punto de Zama neurotóxico y queda \etiqueta{Aturdido}.
\end{CajaEjemplo}

\begin{CajaEjemplo}{Ejemplo 2: Hackeo de un terminal con CRT}
\lead{Situación:} La Biohacker conecta su terminal al armario de distribución de un bloque de viviendas para interceptar la lista de candidatos del Programa Azul.\\
\lead{Ficción:} Suda frío mientras el fósforo verde del monitor parpadea; sabe que está traicionando los protocolos de seguridad que su padre le enseñó a respetar.\\
\lead{Mecánica:} Activa su movimiento «Cacharrería Punk» y tira + \textbf{Culpable} (+1). Saca un doble 5 en los dados: total 11 (\emph{Éxito total}).\\
\lead{Resolución:} Los datos fluyen a su disquete de 3.5 pulgadas. La Biohacker no solo descarga la lista completa del Distrito 7, sino que sobrecarga el conmutador local con un pulso eléctrico que borra sus huellas y deja a oscuras la manzana entera, facilitando su huida.
\end{CajaEjemplo}
